% !TEX root = ../Hauptdatei.tex
% ============================================================
% Kapitel: Listen
% ============================================================

\chapter{Aufzählungen (itemize)}
\label{sec:itemize}
\index{itemize}

\section{Standard-Aufzählung}

Eine einfache Aufzählung mit Standard-Bullet-Punkten:

\begin{itemize}
    \item Erster Punkt
    \item Zweiter Punkt
    \item Dritter Punkt
\end{itemize}

\section{Aufzählung mit eigenem Label}
\index{Label!eigenes}

Mit optionalen Labels lassen sich eigene Symbole verwenden:

\begin{itemize}
    \item[$\star$] Mit Sternchen
    \item[$\rightarrow$] Mit Pfeil
    \item[$\checkmark$] Mit Häkchen
    \item[$\diamond$] Mit Diamant
\end{itemize}

\section{Verschachtelte Aufzählung}
\index{Verschachtelung}

Aufzählungen können verschachtelt werden -- die Symbole ändern sich pro Ebene:

\begin{itemize}
    \item Oberste Ebene
    \begin{itemize}
        \item Zweite Ebene
        \begin{itemize}
            \item Dritte Ebene
        \end{itemize}
        \item Zurück auf Ebene 2
    \end{itemize}
    \item Zurück auf Ebene 1
\end{itemize}

% ============================================================
\chapter{Nummerierungen (enumerate)}
\label{sec:enumerate}
\index{enumerate}
% ============================================================

\section{Standard-Nummerierung}

\begin{enumerate}
    \item Erster Schritt
    \item Zweiter Schritt
    \item Dritter Schritt
\end{enumerate}

\section{Verschiedene Nummerierungsstile}
\index{Nummerierungsstil}

\subsection{Arabische Ziffern (1, 2, 3\ldots):}
\begin{enumerate}[label=\arabic*.]
    \item Eins
    \item Zwei
    \item Drei
\end{enumerate}

\subsection{Kleine Buchstaben (a, b, c\ldots):}
\begin{enumerate}[label=\alph*.]
    \item Alpha
    \item Beta
    \item Gamma
\end{enumerate}

\subsection{Große Buchstaben (A, B, C\ldots):}
\begin{enumerate}[label=\Alph*.]
    \item Erster Buchstabe
    \item Zweiter Buchstabe
    \item Dritter Buchstabe
\end{enumerate}

\subsection{Römische Ziffern (I, II, III\ldots):}
\begin{enumerate}[label=\Roman*.]
    \item Primus
    \item Secundus
    \item Tertius
\end{enumerate}

\section{Nummerierung mit Klammern und Prefix}

\begin{enumerate}[label=\textbf{Aufgabe~\arabic*:}]
    \item Konfiguration laden
    \item Daten validieren
    \item Ergebnis speichern
\end{enumerate}

\section{Verschachtelte Nummerierung mit Referenz}
\index{Verschachtelung}

\begin{enumerate}[label=\arabic*]
    \item Hauptpunkt
    \begin{enumerate}[label=\theenumi.\alph*]
        \item Unterpunkt a
        \item Unterpunkt b
    \end{enumerate}
    \item Hauptpunkt
    \begin{enumerate}[label=\theenumi.\alph*]
        \item Unterpunkt a
        \item Unterpunkt b
        \item Unterpunkt c
    \end{enumerate}
\end{enumerate}

% ============================================================
\chapter{Definitionslisten (description)}
\label{sec:description}
\index{description}
% ============================================================

\section{Standard-Definitionsliste}

\begin{description}
    \item[\LaTeX] Ein Textsatzsystem für professionelle Dokumente
    \item[PDF] Portable Document Format
    \item[Bib\TeX] Ein Literaturverwaltungssystem
\end{description}

\section{Definitionsliste mit Stil-Anpassung}

\begin{description}[style=nextline]
    \item[Itemize] Unnummerierte Aufzählung mit Bullet-Punkten
    \item[Enumerate] Nummerierte Aufzählung mit verschiedenen Zählstilen
    \item[Description] Definitionsliste mit hervorgehobenen Begriffen
\end{description}

% ============================================================
\chapter{Abstand und Formatierung}
\label{sec:abstand}
\index{Abstand}
% ============================================================

\section{Einzug und Abstand anpassen}
\index{Einzug}

Mit \texttt{enumitem}\index{enumitem (Paket)}\cite{mittelbach2004} lassen sich Einzug und vertikaler Abstand steuern:

\begin{itemize}[leftmargin=2cm, itemsep=12pt]
    \item Großer Einzug (2\,cm) und großer Zeilenabstand (12\,pt)
    \item Zweiter Punkt mit gleichem Abstand
    \item Dritter Punkt
\end{itemize}

\section{Kompakte Listen}
\index{Kompaktliste}

Für platzsparende Darstellung:

\begin{itemize}[nosep]
    \item Kein zusätzlicher Abstand zwischen den Punkten
    \item Ideal für Randnotizen oder Fußnoten
    \item Dritte Zeile direkt darunter
\end{itemize}

\section{Fortlaufende Nummerierung}
\index{Fortlaufende Nummerierung}

Zähler über mehrere Listen hinweg weiterführen:

\begin{enumerate}[label=\arabic*., series=beispiel]
    \item Erstes Listenelement
    \item Zweites Listenelement
\end{enumerate}

Hier steht ein Absatz Text zwischen den Listen.

\begin{enumerate}[label=\arabic*., resume=beispiel]
    \item Drittes Listenelement -- Nummerierung wird fortgesetzt
    \item Viertes Listenelement
\end{enumerate}

% ============================================================
\chapter{Farbige Listen}
\label{sec:farbig}
\index{Farbige Listen}
% ============================================================

Mit \texttt{xcolor}\index{xcolor (Paket)} lassen sich einzelne Labels einfärben:

\begin{itemize}
    \item[\textcolor{red}{$\bullet$}] Roter Punkt
    \item[\textcolor{green!60!black}{$\bullet$}] Grüner Punkt
    \item[\textcolor{blue}{$\bullet$}] Blauer Punkt
\end{itemize}

\begin{enumerate}[label=\textcolor{orange}{\textbf{\Alph*.}}]
    \item Orange hervorgehoben
    \item Auch orange
    \item Ebenfalls orange
\end{enumerate}